% Generated deterministically from the frozen Markdown source.
% Responsible author: Carptopus.
\documentclass[10pt]{article}
\usepackage[a4paper,margin=24mm]{geometry}
\usepackage{microtype}
\usepackage{amsmath,amssymb,amsthm,mathtools}
\usepackage{needspace}
\usepackage{xcolor}
\usepackage{xurl}
\usepackage{fontspec}
\newfontfamily\cjkfont{Microsoft YaHei}
\usepackage[colorlinks=true,linkcolor=blue!55!black,citecolor=blue!55!black,urlcolor=blue!65!black]{hyperref}
\usepackage{fancyhdr}

\newenvironment{mdstatement}[1]
  {\par\Needspace{0.16\textheight}\noindent\textbf{#1.}\ \itshape}
  {\par\normalfont}

\setlength{\parindent}{0pt}
\setlength{\parskip}{0.42em}
\setlength{\emergencystretch}{6em}
\setlength{\headheight}{14pt}
\urlstyle{same}

\pagestyle{fancy}
\fancyhf{}
\fancyhead[L]{Carptopus}
\fancyhead[R]{Paving-matroid $h^*$ real-rootedness}
\fancyfoot[C]{\thepage}

\title{Eventual negative real-rootedness for fixed-rank paving matroid base polytopes\\
and a growing-rank root-density theorem}
\author{Carptopus}
\date{7 October 2026}

\hypersetup{
  pdftitle={Eventual negative real-rootedness for fixed-rank paving matroid base polytopes and a growing-rank root-density theorem},
  pdfauthor={Carptopus},
  pdfsubject={Real-rootedness and root density for Ehrhart h-star polynomials of paving matroid base polytopes},
  pdfkeywords={paving matroid, Ehrhart h-star polynomial, real-rootedness, hypersimplex, transfer matrix}
}

\begin{document}
\maketitle
\begin{center}
\fcolorbox{black!35}{black!4}{\parbox{0.90\textwidth}{\centering\small
Version 0.1-beta preprint; not peer reviewed. The full-text comparison with
Adiprasito--Zhang, HAL hal-05747878v1, is recorded in Section 7.}}
\end{center}

\begin{abstract}


Let $M$ be a paving matroid and $h_M^*(x)$ the Ehrhart $h^*$-polynomial
of its base polytope.  We prove two complementary asymptotic results.  First,
for every fixed rank $r\ge3$, all zeros of $h_M^*$ are simple and strictly
negative, uniformly over all sufficiently large $n$-element rank-$r$
paving matroids; no connectivity assumption is imposed.  Second, when the
rank grows with the ground-set size, we give an explicit regime in which a
connected paving matroid has $D-o(D)$ distinct strictly negative zeros,
where $D=\deg h_M^*$.  The second theorem assumes a mild upper bound on the
sizes of effective hyperplanes and does not classify the remaining $o(D)$
zeros.  The proof combines the paving hyperplane-correction formula with
uniform phase grids, endpoint scaling limits, quantitative spectral gaps, and
packing estimates for the hyperplanes.  The two theorems concern different
asymptotic regimes and neither implies the other.

This paper is accompanied by \texttt{proof-\allowbreak{}dossier.\allowbreak{}md} and its frozen 216-file
SHA-256 closure.  The dossier is an integral supplement: whenever a long
uniform contour estimate is summarized in the main text, the cited dossier
section supplies the complete inequalities, endpoint certificates, and
quantifiers.  The theorem statements and assembly arguments do not rely on
historical workflow labels in those records.

\end{abstract}

\section{Main results}


Throughout, $M$ is a loopless rank-$r$ matroid on $[n]$, $P(M)$ is its
base polytope in its intrinsic lattice, and

\[
 1+\sum_{q\ge1}|qP(M)\cap\mathbb Z^n|x^q
 =\frac{h_M^*(x)}{(1-x)^{\dim P(M)+1}}.
\]

A matroid is \textbf{paving} if every circuit has size at least its rank.  An
effective hyperplane is a rank-$(r-1)$ hyperplane $H$ with $|H|\ge r$.
Distinct effective hyperplanes of a paving matroid intersect in at most
$r-2$ elements.
We write $\mathcal H(M)$ for the effective hyperplanes and use the global
convention

\[
 s(M)=\begin{cases}
 \max_{H\in\mathcal H(M)}|H|,&\mathcal H(M)\ne\varnothing,\\
 r-1,&\mathcal H(M)=\varnothing,
 \end{cases}
 \qquad D_u=n-\left\lceil\frac nr\right\rceil. \tag{1.0}
\]

\begin{mdstatement}{Theorem 1.1 (fixed rank)}


For every fixed integer $r\ge3$, there is an integer $N_r$ such that, for
every $n\ge N_r$ and every $n$-element rank-$r$ paving matroid $M$, all
zeros of $h_M^*(x)$ are simple, strictly negative, and real.

\end{mdstatement}

The threshold may depend on $r$; it is not made explicit.  The theorem does
not include rank two or assert anything for all small $n$.

\begin{mdstatement}{Theorem 1.2 (growing rank)}


Let $n_\nu\to\infty$ and $r_\nu\to\infty$.  Put

\[
 D_\nu=n_\nu-\left\lceil\frac{n_\nu}{r_\nu}\right\rceil,
 \qquad X_\nu=\frac{n_\nu}{r_\nu^4},
 \qquad L_\nu=\log X_\nu,
 \qquad A_\nu=\log(2r_\nu+L_\nu).
\]

Assume

\[
 \frac{r_\nu^4A_\nu^4}{L_\nu}\longrightarrow0. \tag{1.1}
\]

Let $M_\nu$ be a connected $n_\nu$-element rank-$r_\nu$ paving matroid,
and suppose every effective hyperplane of $M_\nu$ has size at most
$D_\nu-1$.  Then, for all sufficiently large $\nu$,

\[
 \deg h_{M_\nu}^*=D_\nu
\]

and $h_{M_\nu}^*$ has at least

\[
 D_\nu-J_\nu,
 \qquad J_\nu=\left\lceil r_\nu^7A_\nu^4\right\rceil, \tag{1.2}
\]

pairwise distinct strictly negative real zeros.  Moreover
$J_\nu=o(D_\nu)$.

\end{mdstatement}

We do not assert that the zeros counted in Theorem 1.2 are simple, or that
the remaining $J_\nu$ zeros are real, negative, or distinct.

\section{The exact paving correction formula and the uniform phase grid}


All analytic formulas in Sections 2--6 are first written for a connected
matroid, so that $\dim P(M)=n-1$ and the intrinsic Ehrhart denominator is
$(1-x)^n$.  Disconnected paving matroids are classified and reduced by an
integral affine isomorphism at the end of the fixed-rank proof.

Put $A=1-x$, $m=n-1$, and, for every polynomial $f(k)$, write

\[
 \mathcal E_n(f)=A^n\sum_{k\ge0}f(k)x^k,
 \qquad
 U_{\ell,n}(x)=\mathcal E_n\!\binom{\ell k+n-1}{n-1}. \tag{2.1}
\]

Here $U_{r,n}(x)$ denotes only the leading Veronese term, not the full
uniform-matroid numerator $h_{U_{r,n}}^*(x)$.  The latter also contains
the lower-Euler terms in the first sum of (2.5).
To define the hyperplane correction, let
$M_{\ell,0}(x)$ be the $\ell\times\ell$ matrix, indexed by
$0,\ldots,\ell-1$, whose $(i,j)$-entry is $1$ if $i\ge j$ and $x$
if $i<j$.  Choose any nonzero index $a_0$, let $D_{a_0}(u)$ have entry
$1$ at $a_0$ and $u$ at every other diagonal position, and set

\[
 B_m^{(\ell)}(u,x)
   =e_0^{\mathsf T}\bigl(M_{\ell,0}(x)D_{a_0}(u)\bigr)^m\mathbf 1,
 \qquad
 K_{\ell,m,s}(x)=\sum_{a=0}^{s-1}[u^a]B_m^{(\ell)}(u,x). \tag{2.2}
\]

We abbreviate
$M_\ell(u,x)=M_{\ell,0}(x)D_{a_0}(u)$, and write $M(u)$ when
$\ell,x$ are fixed.
The choice of $a_0\ne0$ does not change the coefficient identity used here.
Indeed,

\[
 [u^a]B_m^{(\ell)}(u,x)
 =A^{m+1}\sum_{k\ge0}
   \binom{(\ell-1)k+a}{a}
   \binom{k+m-a-1}{m-a}x^k. \tag{2.3}
\]

For an effective hyperplane $H$ of size $s$, put $q=n-s$.  For
$0\le h\le r-1$, let $\ell=r-h$ and

\[
 \begin{aligned}
  W_h(k)&=\binom{\ell k+m-h}{m},\\
  Q_{h,H}(k)&=\sum_{J=0}^{k-1}
       \binom{J+q-1}{q-1}
       \binom{\ell k-h-J+s-1}{s-1}.
 \end{aligned} \tag{2.4}
\]

Hanely et al., Proposition 5.10 and Theorem 5.11, give the underlying Ehrhart
polynomial formula.  Multiplying its lattice-point series by $(1-x)^n$,
which is the intrinsic denominator in the connected case, gives the following
full coordinate inclusion--exclusion form:

\[
 h_M^*(x)=
 \sum_{h=0}^{r-1}(-1)^h\binom nh\mathcal E_n(W_h)
 -\sum_H\sum_{h=0}^{r-1}(-1)^h\binom{s_H}{h}
       \mathcal E_n(Q_{h,H}). \tag{2.5}
\]

For $0\le h\le r-2$, the exact finite-band identities are

\[
 \begin{aligned}
  \mathcal E_n(Q_{h,H})
    &=\sum_{j=0}^h(-1)^j\binom hj A^j
      K_{r-h,m-j,s_H-j},\\
  \mathcal E_n(W_h)
    &=\sum_{j=0}^h(-1)^j\binom hj A^j
      U_{r-h,n-j}.
 \end{aligned} \tag{2.6}
\]

When $h=r-1$, one has $Q_{h,H}=W_h$ and
$\mathcal E_n(W_h)=x^{r-1}$; this one-color term is kept separately.
The $h=0$ terms in (2.5) give the exact decomposition

\[
 h_M^*(x)=U_{r,n}(x)-\sum_HK_{r,n-1,s_H}(x)+L_M(x), \tag{2.7}
\]

where $L_M$ is now defined, rather than suppressed, by the terms with
$h\ge1$ in (2.5)--(2.6).  Distinct bad caps are disjoint in the
coordinate-bounded domain because two simultaneous cap violations would force
the total coordinate sum to exceed $rk$; this is the step that makes the sum
over $H$ in (2.5) exact.

We use two packing consequences repeatedly:

\[
 \sum_H\binom{s_H}{r-1}\le\binom n{r-1}, \tag{2.8}
\]

and, for every fixed $r$,

\[
 \sum_Hs_H^{r-1}\le C_r n^{r-1}. \tag{2.9}
\]

In particular, for each fixed $\delta>0$, the number of hyperplanes with
$s_H>\delta n$ is bounded in terms of $r,\delta$, while in the growing-rank
argument it is at most $(C/\delta)^r$.

At $x=-t^r$, the unperturbed matrix in (2.2) is similar to a normal
circulant.  Its eigenvalues are

\[
 \lambda_\omega(t)=\frac{1+t^r}{1-t e^{\pi i/r}\omega},
 \qquad \omega^r=1. \tag{2.10}
\]

Exactly two conjugate branches, denoted $\tau_+$ and $\tau_-$, have
maximal modulus.  We write $S_3$ for the third singular value of the
transfer matrix and set

\[
 d_r(t)=|1-te^{-\pi i/r}|,
 \qquad R_r(t)=\frac{1+t^r}{d_r(t)},
 \qquad
 1-te^{-\pi i/r}=d_r(t)e^{i\phi_r(t)},
 \qquad 0<\phi_r(t)<\frac{r-1}{r}\pi. \tag{2.11}
\]

The two dominant Fourier branches of the leading section $U_{r,n}$ have equal
modulus $R_r(t)^n$.  Their phases lead to the grid

\[
 n\phi_r(t_j)=j\pi. \tag{2.12}
\]

At a grid point, the two principal branches contribute
$2(-1)^jR_r(t_j)^n/r$.  All sign arguments below consist in showing that
the total hyperplane correction and all lower-Euler branches cannot erase
this alternating main term.

The phase has the endpoint expansions

\[
 \phi_r(t)=t\sin(\pi/r)+O_r(t^2)\quad(t\downarrow0), \tag{2.13}
\]

and, with $p=(r-1)/r$,

\[
 p\pi-\phi_r(t)=\frac{\sin(\pi/r)}t
 +\frac{\sin(\pi/r)\cos(\pi/r)}{t^2}
 +O_r(t^{-3})\quad(t\to\infty). \tag{2.14}
\]

The uniform two-branch phase-grid argument follows the hypersimplex prototype
in Adiprasito--Zhang [1, Sections 3.4--3.5].  The estimates below control the
nonuniform paving corrections; the uniform phase mechanism itself is not
claimed as new.

\section{Quantitative sign control}


This section packages the analytic estimates used in the growing-rank proof.
They are also the local ingredients of the fixed-rank argument.

\begin{mdstatement}{Proposition 3.0 (the ultra-small direct endpoint)}


Let $r,n\to\infty$, put $m=n-1$, and let $t_j\le1/4$ be a phase-grid
point.  Set $\kappa=mt_j$.  Suppose, with fixed absolute constants
$C_*,C_1$, that

\[
 \begin{gathered}
 \frac{\kappa}{r^8\log^2(2r)}\longrightarrow\infty,
 \qquad
 \frac{\kappa}{r^6\log^6(1+\kappa)}\longrightarrow\infty,\\
 r^2+r\log(2r)=o(\log(1/t_j)),\\
 C_*^rt_j\log(1+\kappa)=o(1/r),\qquad
 (C_1r)^rC_*^rt_j\log(1+\kappa)=o(1/r).
 \end{gathered} \tag{3.0}
\]

If every effective hyperplane has size at most
$D-1=n-\lceil n/r\rceil-1$, then

\[
 (-1)^j\frac{h_M^*(-t_j^r)}{R_r(t_j)^n}\ge\frac1{2r}. \tag{3.0a}
\]

\end{mdstatement}

\begin{proof}


Take $\delta=1/[1000(r+1)]$.  The paving packing inequality gives at most
$(C_1r)^r$ hyperplanes with $s_H>\delta n$, while
$\sum_H(s_H)_{r-1}\le4^{r-1}n^{r-1}$.  The small hyperplanes and every
lower-Euler layer in (2.5)--(2.6) have total normalized size

\[
 \exp\{C r^2\log(2r)\}(1+\kappa)^{2r}e^{-c\kappa}=o(1/r). \tag{3.0b}
\]

All but at most one large hyperplane are separated from the critical density.
For each such term, the central tail, middle arc, and far arc give

\[
 C_*^rt_j\log(1+\kappa)
 +P_r(\kappa)e^{-c\sqrt\kappa}
 +P_r(\kappa)e^{-c\kappa/r^2+C\sqrt\kappa}
 +P_r(\kappa)e^{-c\kappa/(r^4\log(2r))+C\sqrt\kappa}, \tag{3.0c}
\]

where $\deg P_r\le Cr$ and its coefficient height is
$\exp(O(r^2\log(2r)))$.  The five assumptions in (3.0), including the last
one after multiplication by $(C_1r)^r$, make their full sum $o(1/r)$.
For the unique possible critical hyperplane, the complex Gaussian cumulative
has modulus at most $1/2$, so its signed contribution is at most
$1/r+o(1/r)$.  The uniform pair contributes $2/r$, proving (3.0a).

\end{proof}


\begin{mdstatement}{Proposition 3.1 (central and noncentral hyperplane estimates)}


Let $r\ge4$, $m=n-1$, and let $s\le (r-1)m/r$.  Put

\[
 \xi=\frac{s-(r-1)m/r}{\sqrt m}.
\]

On the phase grid (2.12), the normalized highest-Euler correction has the
central approximation

\[
 \frac{K_{r,m,s}(-t^r)}{R_r(t)^n}
 =\frac2r\operatorname{Re}\left[
 e^{in\phi_r(t)}
 \Phi\left(\frac{\xi}{\sqrt{v_r(t)}}\right)
 \right]+E_{r,n,t,s}, \tag{3.1}
\]

where

\[
 v_r(t)=\frac{(r-1)/r}{rt\,e^{\pi i/r}}, \tag{3.2}
\]

and the branch of the Gaussian distribution function satisfies

\[
 \left|\Phi\left(\frac{\xi}{\sqrt{v_r(t)}}\right)\right|\le\frac12.
 \tag{3.3}
\]

The error in (3.1) is $o(1/r)$, uniformly in the following three regions,
under the corresponding assumptions:

\begin{enumerate}
\item $t_-\le t\le1/4$, where
$t_-=n^{-1/(K_0r)}$ and $r^3\log(2r)=o(\log n)$;
\item $1/4\le t\le4$, where
\[
e^{C_0r\log(2r)}r^{12}\log^3(n+3)=o(n); \tag{3.4}
\]
\item $t\ge4$, with $\kappa=(n-1)/t$, where
\[
\frac{\kappa}{r^6\log^3(2r+\kappa)}\longrightarrow\infty
\tag{3.5}
\]
and (3.4) holds.
\end{enumerate}

For hyperplanes separated from the critical ratio $(r-1)/r$, the right
side is $o(1/r)$ after summation over all such hyperplanes.  At most one
effective hyperplane can lie in the critical window.

\end{mdstatement}

\begin{proof}


The content transfer matrix at $x=-t^r$ has two maximal-modulus Fourier
eigenvalues and $r-2$ subordinate eigenvalues.  On $1/4\le t\le4$, the
third singular value $S_3$ and the two dominant phases satisfy

\[
 1-(S_3/R_r)^2\ge c/r^2,
 \qquad |\tau_+-\tau_-|\ge c/r. \tag{3.6}
\]

If $u=e^{i\theta}$ is the coefficient-extraction variable, the rank-one
diagonal perturbation and projection onto the dominant two-dimensional
Fourier space give

\[
 \frac{\operatorname{spr}M(u)}{R_r}
 \le1-c\min\{\theta^2/r^3,r^{-4}\}. \tag{3.7}
\]

Schur triangularization then yields

\[
 \left\|\left(M(u)/R_r\right)^m\right\|
 \le\exp\{Cr\log(m+r)-cm\min(\theta^2/r^3,r^{-4})\}. \tag{3.8}
\]

In a $c/r^2$ neighborhood of $u=1$, disjoint root circles and the
resolvent formula bound the first three logarithmic derivatives of the two
principal branches and their residual amplitudes by

\[
 Q_r\le(Cr)^{Cr}. \tag{3.9}
\]

Using the inward contour
$u=\exp((-h+iy)/\sqrt m)$, completing the quadratic exponent, and using
(3.8) off the central arc gives (3.1).  The argument of
$-\xi/\sqrt{2v_r(t)}$ lies in the positive-kernel sector, which proves
(3.3).

For $0<t\le1/4$, the corresponding estimates are

\[
 1-(S_3/R_r)^2\ge ct/r^2,
 \qquad |\tau_+-\tau_-|\ge ct/r, \tag{3.10}
\]

and hence

\[
 \operatorname{spr}(M(e^{i\theta}))/R_r
 \le1-c\min\{t\theta^2/r^3,t^3/r^4\}. \tag{3.11}
\]

The analytic envelope is at most $(Cr/t)^{Cr}$.  Taking
$t\ge n^{-1/(K_0r)}$ with $K_0$ sufficiently large makes the central
Taylor error and the outer-arc loss uniform, proving the first region.

For $t\ge4$, scale the characteristic polynomial by
$\lambda=t^{r-1}\mu$, $\varepsilon=1/t$.  It becomes

\[
 Q_{r,t}(\mu,u)
 =(\varepsilon\mu-a_0)(\varepsilon\mu-ua_0)^{r-1}+\mu^r,
 \qquad a_0=1+t^{-r}. \tag{3.12}
\]

Its limiting roots satisfy
$\nu_j^r=(-1)^{r+1}u^{r-1}$.  Rouché circles of radius $c/r$ and one
step of perturbation theory give

\[
 \mu_j=\nu_j-\frac1t\frac{\nu_j^2}{r}
 \left(1+\frac{r-1}{u}\right)+O(Q_r/t^2), \tag{3.13}
\]

with the same bound after three logarithmic derivatives.  Consequently the
principal modulus loss is at least $c\theta^2/(rt)$ near zero and
$c/(tr^3)$ away from zero.  The contour scale
$r\sqrt{B\log(2r+\kappa)}/\sqrt\kappa$, for a fixed profile constant
$B$, and (3.5) now give (3.1).

Finally, (2.8) implies that hyperplanes above a fixed density have cardinality
at most $(Cr)^r$.  Two hyperplanes cannot both meet the critical density
window, because their intersection would exceed $r-2$.  The retained
Gaussian left-tail factor for every noncritical term absorbs this count.
Small hyperplanes are summed before normalization, using (2.8), rather than by
multiplying a pointwise $o(1)$ estimate.  This proves the proposition.

\end{proof}


\begin{mdstatement}{Proposition 3.2 (low-Euler suppression)}


Under the hypotheses of each of the three regions in Proposition 3.1,

\[
 \frac{L_M(-t^r)}{R_r(t)^n}=o(1/r) \tag{3.14}
\]

uniformly over the indicated paving matroids.

\end{mdstatement}

\begin{proof}


The layer of Euler index $r-h$ has spectral-radius ratio

\[
 \frac{R_{r-h}}{R_r}\le
 \begin{cases}
  \exp(-ct/r^3),&0<t\le1/4,\\
  \exp(-c/r^3),&1/4\le t\le4,\\
  \exp(-c/(tr^3)),&4\le t\le T_r,\\
  \exp(-c\log t/r),&t\ge T_r,
 \end{cases} \tag{3.15}
\]

for an exponential threshold $T_r$.  The finite-band, inclusion--exclusion,
and hyperplane-packing factors contribute at most

\[
 \exp\{Cr\log n+Cr^2\log(2r)+Cr^2\log^+t\}. \tag{3.16}
\]

The assumptions in the respective regions make the $n$-fold loss from
(3.15) dominate (3.16).  The isolated one-color term has the stronger bound
$\exp(Cr\log n+Cr^2\log t-cnr\log t)$ in the large-$t$ region.
Thus all lower-Euler layers sum to (3.14).
\end{proof}


\begin{mdstatement}{Corollary 3.3 (alternating grid signs)}


Whenever the hypotheses of Propositions 3.1 and 3.2 hold at a phase-grid
point $t_j$,

\[
 (-1)^j\frac{h_M^*(-t_j^r)}{R_r(t_j)^n}\ge\frac1{2r} \tag{3.17}
\]

for all sufficiently large parameters.

\end{mdstatement}

\begin{proof}


The uniform pair contributes $2/r$.  The unique possible critical
hyperplane contributes at most $1/r+o(1/r)$, while all noncritical
hyperplanes and all lower-Euler layers contribute $o(1/r)$.
\end{proof}


\begin{mdstatement}{Proposition 3.3A (bounded critical-excess middle signs)}


Fix $r\ge4$ and $C\ge0$.  Put

\[
 D_u=n-\left\lceil\frac nr\right\rceil,
 \qquad s=D_u+\delta,
 \qquad 0\le\delta\le C. \tag{3.17a}
\]

For all sufficiently large $n$, every connected rank-$r$ paving matroid
with largest effective hyperplane of size $s$ has the alternating sign
(3.17) at every phase-grid point between the fixed direct and reciprocal
endpoint circles.  The threshold is uniform in the matroid and in
$\delta\in\{0,\ldots,C\}$.

\end{mdstatement}

\begin{proof}


Write

\[
 \begin{aligned}
 b_u&=r\lceil n/r\rceil-n,\\
 \ell&=\left\lceil\frac{\delta+1-b_u}{r-1}\right\rceil,\\
 c&=\lceil n/r\rceil+\ell,\qquad b=b_u+r\ell,
 \end{aligned}
 \qquad
 a=(r-1)\ell+b_u-\delta-1. \tag{3.17b}
\]

Thus $a\in\{0,\ldots,r-2\}$, while $a,b,\ell$ range over a finite set
depending only on $r,C$.  Insert the shifted highest cutoff into the exact
finite-band formula (2.5)--(2.7).  Relative to the central cutoff in
Proposition 3.1, the two-moving local expansion is multiplied by
$\exp(O(\sigma))=1+o(1)$, uniformly because the shift is bounded.  The
central Gaussian parameter still satisfies $\xi\to0$.  The packing estimate
for every additional hyperplane and the complete lower-Euler estimate are
unchanged, and hence their sum is $o(1/r)$.  The uniform pair therefore
contributes $2/r+o(1/r)$, the unique possible critical hyperplane at most
$1/r+o(1/r)$, and all remaining terms $o(1/r)$.  This yields the strict
margin $1/(2r)$ on the whole middle grid.

This is an independent bounded-increment estimate, not an application of
Proposition 3.1 outside its stated range.  The complete Taylor bounds,
including the two-moving remainder and lower-Euler sum, are in proof-dossier
D2 under S59 and S91.  Since only finitely many $(a,b,\ell)$ occur, their
thresholds have a common maximum.
\end{proof}


\begin{mdstatement}{Proposition 3.4 (the giant-profile middle sign)}


Fix $r\ge5$, put $d=r-1$, and suppose the largest effective hyperplane is
$H_0$ of size $n-q$, where $2\le q<n/r$.  Set
$p=q/n$, $a_0=1-rp$, and let
$W=h e^{i\theta}$, $\pi/r<\theta<\pi/d$, be the leading solution of

\[
 (1-p)W^r+a_0t^rW+dp\,t^r=0. \tag{3.18}
\]

In the separated subcritical range $p\le p_{\max}<1/r$, define

\[
 \chi_n(t)=n[-\arg(1-W)+p\arg(dp+a_0W)]
 -\frac\theta2-\frac12\arg(rp+a_0W). \tag{3.19}
\]

There is a fixed-rank threshold $L=L(r,p_{\max})$ such that, whenever

\[
 q,\qquad n-rq,\qquad nh,\qquad n/h\ge L, \tag{3.20}
\]

every point satisfying $\chi_n(t)=\ell\pi$ has
$\operatorname{sign}h_M^*(-t^r)=(-1)^\ell$.  The statement includes every
legal collection of additional hyperplanes.

In the moving critical range $p\to1/r$, $n-rq\to\infty$, put

\[
 U_n(t)=n[-\arg(1-W)+p\arg(dp+a_0W)]. \tag{3.21}
\]

There is a threshold $L_r$ such that the same four inequalities (3.20) and
$U_n(t)=\ell\pi$ imply the same sign.  No hypothesis that
$n(1-rp)^2\to\infty$ is required.

\end{mdstatement}

\begin{proof}


Put

\[
 K=\frac{a_0h}{dp},\qquad
 V=p(1-p)\left(a_0+\frac{rp}{W}\right),
 \qquad \Sigma=K\sqrt{n\operatorname{Re}V}. \tag{3.21a}
\]

The exact contour obtained from (2.5)--(2.7) is split into the two leading
saddle basins, their bounded connecting arc, and the remaining arcs.  The
full estimates in proof-dossier D4 give, uniformly over every legal family of
additional hyperplanes,

\[
 h_M^*(-t^r)=2\mathcal M\cos\chi_n+E,\qquad
 \frac{|E|}{\mathcal M}\le
 \varepsilon_r(q,n-rq,nh,n/h,\Sigma), \tag{3.21b}
\]

where $\mathcal M>0$, and $\varepsilon_r\to0$ whenever its five displayed
arguments tend to infinity.  This is a full-contour estimate: the central arc
uses the quadratic spectral gap, the bounded connecting arc uses the exact
cumulative integral, and the far arcs use the strict full-circle minimum-modulus
gap.  Packing is applied before summing the additional-hyperplane corrections,
so no factor equal to their number remains.  In the subcritical range,
$\Sigma^2\ge c_{r,p_{\max}}nh$; hence the four inequalities (3.20) make
the right side of (3.21b) smaller than the fixed positive grid margin.  This
proves the first clause.

For the moving critical range, a sequence violating the statement has a
subsequence on which $\Sigma$ either diverges or stays bounded.  The first
case again follows from (3.21b).  In the bounded case, the exact moving-saddle
relations imply $a_0\to0$, $h/t\to1$,
$na_0^2t=O(1)$, and $a_0(1+t)\to0$.  Taylor expansion of the Legendre
exponent, uniformly on the real parameter segment, gives

\[
 n\{S(p)-S(1/r)\}
 =-\frac{n(1-rp)^2W_0}{2d}
  +O\!\left(n(1-rp)^3t(1+t)\right),
 \qquad W_0=t e^{\pi i/r}. \tag{3.22}
\]

With

\[
 z=-\frac{\xi}{\sqrt{2v_r(t)}},\qquad
 \xi=-\frac{na_0+d}{r}\sqrt{\frac{t}{n-1}}, \tag{3.22a}
\]

this is equivalent to

\[
 U_n(t)=n\phi_r(t)-\operatorname{Im}z^2+o(1). \tag{3.22b}
\]

The remaining scaled complementary-error factor is
$E(z)=e^{z^2}\operatorname{erfc}(z)$, with
$\arg z=\pi/(2r)$.  Its exact integral representation

\[
 E(z)=\frac2\pi\int_0^\infty
 \frac{z e^{-u^2}}{u^2+z^2}\,du \tag{3.22c}
\]

puts every integrand in the sector
$[ -\pi/(2r),\pi/(2r)]$.  Thus
$\operatorname{Re}E(z)\ge\cos(\pi/(2r))|E(z)|>0$, uniformly on each bounded
separation window.  At $U_n=\ell\pi$, (3.22b)--(3.22c) therefore give sign
$(-1)^\ell$.  Divergent and bounded $\Sigma$ exhaust every failure
sequence, so the sequential estimate is equivalent to one finite threshold
$L_r$ for the four quantities in (3.20).
\end{proof}


\begin{mdstatement}{Proposition 3.5 (rank-four giant-profile middle signs)}


Let $r=4$, $d=3$, and let the largest effective hyperplane have size
$n-q$.  Define $p=q/n$, $a_0=1-4p$, the leading solution
$W=h e^{i\theta}$ by (3.18), and the phases $\chi_n$ and $U_n$ by
(3.19) and (3.21).  There are finite gates with the following properties.

\begin{enumerate}
\item For each $p_{\max}<1/4$, if $q$, $nh$, and $n/h$ exceed the gate,
then every $\chi_n(t)=\ell\pi$ between the endpoint circles has sign
$(-1)^\ell$.
\item If $q$, $n-4q$, $nh$, and $n/h$ exceed the critical gate, then
every $U_n(t)=\ell\pi$ between the endpoint circles has sign
$(-1)^\ell$.
\end{enumerate}

Consequently these signs hold for every actual connected sequence in one of
the three ranges

\[
 q\to\infty,\ q/n\to0,\qquad q/n\to p_0\in(0,1/4),\qquad
 q/n\to1/4\text{ with }n-4q\to\infty, \tag{3.23}
\]

uniformly over all additional hyperplanes.

\end{mdstatement}

\begin{proof}


For the separated regime, the exact four-state transfer matrix has a simple
leading conjugate pair on the full contour.  After normalizing by its leading
modulus, the other two eigenvalues and every attached-hyperplane correction
satisfy

\[
 \frac{|E|}{\mathcal M}
 \le \varepsilon_4(q,n-4q,nh,n/h), \tag{3.24}
\]

where the right side tends to zero when the displayed active arguments do.
This estimate is obtained by the same basin/connecting-arc/far-arc partition
as (3.21b), but with the quartic characteristic polynomial retained exactly;
the full-circle spectral gap and the resolvent bounds are in proof-dossier D5,
S50--S54.  At a $\chi_n$-grid point the normalized leading pair is
$2\cos\chi_n$, so (3.24) proves the first clause.

When $p\to1/4$, write again
$z=-\xi/\sqrt{2v_4(t)}$.  The exact rank-four Legendre calculation gives

\[
 U_n(t)=n\phi_4(t)-\operatorname{Im}z^2+o(1), \tag{3.25}
\]

uniformly in the bounded separation window; the remainder is controlled by
the four quantities in the second clause, without an additional $\Sigma$
gate.  The scaled factor is $e^{z^2}\operatorname{erfc}z$, and the sector
argument (3.22c) gives its strictly positive real part.  This proves the
critical clause.  If $q/n\to0$ with $q\to\infty$, use the separated
clause after the sparse-outside rescaling; bounded $q$ is handled separately
by Proposition 4.4.  The paving intersection bound is imposed before summing
the exact corrections, so (3.24)--(3.25) concern the complete polynomial.

\end{proof}


\section{Endpoint root blocks}


The sign grid does not by itself reach either endpoint.  We use reciprocal
scalings and Rouché transfer.

\begin{mdstatement}{Proposition 4.0 (fixed-rank direct endpoint)}


Fix $r\ge4$, put $p=(r-1)/r$, and let
$\boldsymbol\alpha=(\alpha_1,\alpha_2,\ldots)$ be a finite or countable
sequence with $\alpha_i\ge0$, $\sum_i\alpha_i\le1$, and
$\max_i\alpha_i\le p$.  If
$(N_1,N_2,\ldots,N_0)$ has multinomial distribution with $rj$ trials and
cell probabilities $(\alpha_1,\alpha_2,\ldots,1-\sum_i\alpha_i)$, define

\[
 F_{\boldsymbol\alpha}(z)=\sum_{j\ge0}
 \Pr\{N_i\le(r-1)j\text{ for every }i\}\frac{z^j}{(rj)!}. \tag{4.0}
\]

Every $F_{\boldsymbol\alpha}$ has infinitely many simple strictly negative
zeros.  For every sufficiently large fixed $j$, the circle

\[
 |z|=T_j^r,\qquad T_j=\frac{\pi j}{\sin(\pi/r)}, \tag{4.0a}
\]

contains exactly $j$ zeros and has no boundary zero, uniformly over all
such spectra.  If $M_n$ is any connected fixed-rank paving sequence whose
effective-hyperplane densities converge to $\boldsymbol\alpha$, then
$h_{M_n}^*(z/n^r)\to F_{\boldsymbol\alpha}(z)$ locally uniformly.  Hence,
uniformly over the actual objects in Lemma 5.1, the interval
$0<t<t_{n,j}$, where $n\phi_r(t_{n,j})=j\pi$, contains exactly $j$
simple negative roots and the boundary value has sign $(-1)^j$.

\end{mdstatement}

\begin{proof}


Formula (2.5) gives the coefficient probability in (4.0), and factorial
domination gives compact convergence, including countable spectra.  We first
prove the root statement without using the circles.  Put $d=r-1$, order the
positive $\alpha_i$, and set

\[
 h(y)=\frac{y}{ry-d},\qquad f(t)=\frac{t}{d-rt},
 \qquad N(t)=\#\{i:\alpha_i>t\}. \tag{4.0b}
\]

Root-of-unity filtering and the common two-branch change of variables transform
the normalized coefficient series into a cosine transform with signed weight
$H$.  Its downward variation is exactly

\[
 V_-=\sum_{i\ge2}f(\alpha_i)-h(p_b)+1,
 \qquad b=\alpha_2, \tag{4.0c}
\]

with the empty or one-atom cases obtained by deletion.  Here
$p_b\in(d/r,1)$ is the unique high-branch point paired with the
low-branch point $b$, namely
$b^d(1-b)=p_b^d(1-p_b)$.  Thus $h(p_b)>1$ for $b>0$.
Since $\sum_{i\ge2}\alpha_i\le1-b$ and $\alpha_i\le b\le1/2$,

\[
 V_-<\frac{1-b}{d-rb}\le\frac1{r-2}\le\frac12. \tag{4.0d}
\]

All jumps of $H$ are upward, $H(0)=1$, and the possible terminal
inverse-square-root singularity is integrable.  Truncating that singularity
does not increase $V_-$.  Hence the signed-cosine rectangle argument has
strict vertical and horizontal margins and puts the normalized entire function
$C$ in the Laguerre--Pólya class.  The exact gamma multiplier

\[
 \gamma_j=\frac{(2j)!}{(dj)!j!} \tag{4.0e}
\]

takes $C$ to $F_{\boldsymbol\alpha}$; its multiplier-sequence
strictification makes every zero simple and negative.  The coefficient is
strictly positive, the order is at most $1/r<1$, and the function is not a
polynomial, so there are infinitely many zeros.  The countable case follows
by truncation: $tN(t)\to0$, the total jump mass is finite, and the kernel
integrals are dominated by $\sum_i\alpha_i^{1+d/2}<\infty$.

Now all zeros are already known to be simple and negative.  The explicit
root-of-unity formula for the unperturbed kernel gives the sign $(-1)^j$ at
$-T_j^r$.  The same rectangle margins and (4.0d) make that sign uniform over
the compact spectrum space for all sufficiently large fixed $j$.  Thus the
circle $|z|=T_j^r$ is zero-free: any boundary zero would have to be the
negative point already excluded by the strict sign.  The negative-axis sign
changes provide $j$ real zeros inside, and the argument principle gives the
same count.  Local uniform convergence of actual rescaled polynomials and
small disjoint Rouché disks transfer the count and simplicity.  A failure
sequence has a convergent spectrum subsequence and contradicts this zero-free
circle.  The full cosine-transform calculation and the uniform rectangle
estimates are proof-dossier D4.
\end{proof}


\begin{mdstatement}{Proposition 4.1 (fixed-rank endpoint transfer)}


Fix $r\ge4$, let $b=r\lceil n/r\rceil-n$, and suppose first that the
largest effective hyperplane has size $s\le D_u-1$, where
$D_u=n-\lceil n/r\rceil$.  Put $g=D_u-s\ge1$.  The reciprocal endpoint
kernel is

\[
 F_{r,b,g}(z)=\sum_{j\ge0}
 \Pr\!\left\{\operatorname{Bin}
   \left(rj+b,\frac{r-1}{r}\right)
   \le(r-1)j+b+g-1\right\}\frac{z^j}{(rj+b)!}. \tag{4.1}
\]

For every $b\in\{0,\ldots,r-1\}$ and $g\ge1$, this entire function has
infinitely many simple strictly negative zeros.  As $g\to\infty$, it
converges locally uniformly to

\[
 V_{r,b}(z)=\sum_{j\ge0}\frac{z^j}{(rj+b)!}, \tag{4.2}
\]

which has the same zero property.  This limiting function is the residue
function $E_b$ of Adiprasito--Zhang [1, Section 3.3], with $k=r$; its
simple negative zeros are prior work.  Here the additional assertions concern
the corrected kernels $F_{r,b,g}$ and counts uniform in $g$.
For every sufficiently large fixed $k$,
the circles

\[
 |z|=T_{k,b}^r,
 \qquad T_{k,b}=\frac{\pi(k+b/r)}{\sin(\pi/r)}, \tag{4.3}
\]

contain exactly $k$ zeros, have no boundary zero, and work uniformly in
$g$, including the limit $g=\infty$.

If instead $0\le s-D_u\le C$, the actual codegree differs from
$\lceil n/r\rceil$ by a bounded integer.  For fixed $r,C$, the residue
$b$ and the additional endpoint index $a$ range over a finite set.  Denote
by $E_{r;b,a}$ the entire function obtained as the locally uniform limit of
the exact reciprocal polynomial (2.5) in that fixed $(b,a)$-subsequence.
Every such $E_{r;b,a}$ has infinitely many simple strictly negative zeros
and zero-free counting circles of the form (4.3).  Consequently every fixed
finite block of direct and reciprocal endpoint zeros transfers uniformly to
simple negative zeros of $h_M^*$.

\end{mdstatement}

\begin{proof}


Coefficientwise convergence follows from (2.3)--(2.7), and uniform factorial
tails make it locally uniform at the origin.  The sign of
$F_{r,b,g}(-T_{k,b}^r)$ is $(-1)^k$.  A full-circle argument-principle
count gives exactly $k$ zeros inside (4.3); hence the alternating real zeros
already supplied by the signs exhaust the disk and are simple.  The estimates
are uniform in $g$: if a common transfer threshold failed, a failure sequence
would have, after fixing the finite residue $b$, either a constant subsequence
$g=g_0$ or a subsequence $g\to\infty$.  The corresponding locally uniform
limit is respectively $F_{r,b,g_0}$ or $V_{r,b}$, and the same zero-free
circle contradicts failure by Rouché's theorem.  Thus no maximum over infinitely
many individual $g$-thresholds is taken.  The bounded-excess kernels form a
genuinely finite $(b,a)$-family, so the same argument takes a finite maximum
there.
\end{proof}


\begin{mdstatement}{Proposition 4.2 (growing-rank reciprocal block)}


Under (1.1), let

\[
 K=\left\lceil r^5A^4\right\rceil,
 \qquad K_{\max}=\left\lfloor\frac r{16}L\right\rfloor. \tag{4.4}
\]

Set $b=r\lceil n/r\rceil-n$ and define

\[
 Q_n(z)=n^{-b}\left(\frac z{n^r}\right)^D h_M^*(n^r/z). \tag{4.5}
\]

Then $K/K_{\max}\to0$, and $Q_n$ has exactly $K-1$
simple strictly negative zeros in $|z|<T_{K-1,b}^r$, with no zero on the
boundary and boundary sign $(-1)^{K-1}$.  Under (4.5), these correspond to
$K-1$ zeros of $h_M^*$ in
$x<-n^r/T_{K-1,b}^r$.

\end{mdstatement}

\begin{proof}


For the first $2K$ reciprocal coefficients, the paving correction formula
has the same leading factorial kernel as the uniform polynomial, with a
relative error tending to zero uniformly in the hyperplane family.  On the
circle of radius $T_{K-1,b}^r$, with $T_{k,b}$ as in (4.3), the coefficient
tail is split at the central
index.  Below it, the factorial ratio is geometric after the $K$-scale;
above it, the reciprocal factorial tail is dominated by its first omitted
term.  Condition (1.1) implies $K=o(rL)$, so the two estimates overlap.
More precisely, if

\[
 G_{r,b,K}(z)=\sum_{j=0}^{2K}\frac{z^j}{(rj+b)!},
 \qquad \Gamma_K=\{|z|=T_{K-1,b}^r\}, \tag{4.5a}
\]

then the two coefficient ranges and their tails give

\[
 \sup_{z\in\Gamma_K}|Q_n(z)-G_{r,b,K}(z)|
 \le \eta_n
 \inf_{z\in\Gamma_K}|G_{r,b,K}(z)|,
 \qquad \eta_n\longrightarrow0. \tag{4.5b}
\]

The denominator in (4.5b) is positive by the full root-of-unity circle
estimate, not by a negative-axis test.  The coefficient-ratio derivation of
(4.5b), including the uniform paving correction and both factorial tails, is
proof-dossier D6.  Rouché's theorem now gives the exact root count and boundary
sign for $Q_n$.
Multiplication by the monomial used in (4.5) changes the boundary sign by the
known parity $(-1)^D$, which is the conversion used in (6.7).
\end{proof}


\begin{mdstatement}{Proposition 4.3 (fixed-rank giant-profile endpoint blocks)}


Fix $r\ge4$, put $d=r-1$, and suppose the largest effective hyperplane
has size $n-q$, where $2\le q<n/r$.  With
$p=q/n$, $a_0=1-rp$, and the leading solution $W=h e^{i\theta}$ of
(3.18), put

\[
 c=\left\lceil\frac{n-q+1}{r-1}\right\rceil,\qquad
 b=rc-n,\qquad a=(r-1)c-(n-q)-1,\qquad D=n-c. \tag{4.6}
\]

Then $a\in\{0,\ldots,r-2\}$, the codegree is exactly $c$, and the degree
is exactly $D$.  In the subcritical range $q\to\infty$,
$q/n\le p_{\max}<1/r$, every sufficiently large fixed $j,k$ admits
disjoint boundaries

\[
 x_L=-\frac{Y_j}{p(1-p)^{r-1}n^r},\qquad
 x_U=-\frac{\Lambda}{T_{a,k}^{r-1}},\qquad
 \Lambda=\frac{(n-q)^{r-1}(q+b)}b, \tag{4.7}
\]

such that the disk ending at $x_L$ contains exactly $j$ simple negative
roots and has sign $(-1)^j$, while the exterior beyond $x_U$ contains
exactly $k$ simple negative roots and has sign $(-1)^{D-k}$.  Here
$Y_j$ is the fixed Wright-phase radius with its phase equal to $j\pi$, and
$T_{a,k}=\pi(k+a/(r-1))/\sin(\pi/(r-1))$.

In the critical range $q/n\to1/r$, $n-rq\to\infty$, the same conclusion
holds with the direct boundary
$x_L=-R_{p,j}/n^r$, where the fixed macroscopic circle $R_{p,j}$ contains
exactly $j$ roots, and with the same reciprocal boundary (4.7).  At these
boundaries,

\[
 U_n(t_L)=j\pi+o(1),\qquad
 U_n(t_U)=(D-k)\pi+\frac{\pi}{r-1}+o(1). \tag{4.8}
\]

\end{mdstatement}

\begin{proof}


In the subcritical range, multiplication of the exact saddle equation by
$n^r$ at the scale
$t^r=Y/[p(1-p)^{r-1}n^r]$ gives

\[
 nW\longrightarrow
 \frac{d^{1/r}Y^{1/r}}{1-p_0}e^{\pi i/r}. \tag{4.8a}
\]

The exact coefficient formula and factorial domination therefore give local
uniform convergence of the direct polynomial to the Wright kernel whose
continuous phase $\Psi$ satisfies $\Psi(Y_j)=j\pi$.  Its full Hankel-circle
estimate has one strictly dominant conjugate pair on the boundary and hence
exactly $j$ simple negative zeros inside.  The remaining macroscopic
hyperplanes contribute an integrable signed weight whose downward variation is
strictly below one half; thus the same boundary is zero-free for the complete
kernel.  The convergence is uniform in all legal spectra, so Rouché transfers
the count to the actual polynomial.

At the reciprocal scale, after fixing
$a\in\{0,\ldots,d-1\}$, the normalized polynomial converges locally
uniformly to

\[
 V_{d,a}(z)=\sum_{m\ge0}\frac{z^m}{(dm+a)!}. \tag{4.8b}
\]

The function (4.8b) is likewise the residue function of [1, Section 3.3],
with $k=d$.  We use its known simple negative zeros as the uniform endpoint
baseline, rather than claim a new zero theorem for that function.
The root-of-unity filter shows that the circle
$|z|=T_{a,k}^{d}$ has one dominant conjugate pair and is zero-free, with
exactly $k$ simple negative zeros inside.  The complete correction is
uniformly smaller on that whole circle, not merely at its negative point.
Rouché therefore gives the reciprocal count and sign in (4.7).

In the critical range, the direct limit is the macroscopic Beta kernel.  Its
signed-weight representation has positive initial mass, upward jumps, and
negative variation below one half, so the same cosine-transform rectangle
argument gives only simple negative zeros and zero-free counting circles.  The
reciprocal limit remains (4.8b).  The local-uniform transfers at both circles
are uniform when $q,n-rq,nh,n/h$ exceed their finite gate.  The explicit
Hankel, Beta, and reciprocal boundary estimates are proof-dossier D4.

For the degree, every interior point must have dilation
at least both $\lceil n/r\rceil$ and $\lceil(n-q+1)/(r-1)\rceil=c$.
At dilation $c$, put one in every coordinate and distribute the additional
$b$ units among the $q$ coordinates outside $H_0$, never exceeding
$c-2$ extra units per coordinate.  The capacity inequality
$q(c-2)\ge b$ holds for large $n$.  Then $x(H_0)=n-q<(r-1)c$; every
other effective hyperplane has size at most $q+r-2$, and hence has sum at
most $q+r-2+b<(r-1)c$.  This is a relative interior lattice point, proving
the codegree and degree assertions.  The phase labels in (4.8) follow by
expanding the exact saddle equation (3.18); notably the finite
$\pi/(r-1)$-offset is retained.
\end{proof}


\begin{mdstatement}{Proposition 4.4 (fixed-complement theorem)}


Fix integers $r\ge3$ and $q\ge2$.  For all sufficiently large $n$,
every connected rank-$r$ paving matroid having an effective hyperplane of
size $n-q$ has only simple strictly negative $h^*$-zeros, with no
restriction on the remaining legal effective hyperplanes.

\end{mdstatement}

\begin{proof}


Put $k=r-1$, $m=q-1$, and $N=n-q+1$.  For the one-hyperplane model,
split the outside-coordinate sum into the intervals
$s=(b+1)t+y$, $0\le y\le t$, and inclusion--exclude the upper bounds.
The outside fibre is

\[
 F_b(t,y)=\sum_{v=0}^{b+1}(-1)^v\binom qv
 \binom{(b+1-v)t+y+q-1-v}{q-1}. \tag{4.9}
\]

Expanding in $\binom yj$, the required uniform-slice moments are

\[
\begin{aligned}
 M_{a,N,j}(t)
 &=\sum_{i=0}^{a-1}(-1)^i\binom{N-1}{i}
   \binom{(a-i)t+N-1-i}{N+j-1}\\
 &\quad-
 \sum_{i=0}^{a-2}\sum_{\ell=0}^{j}(-1)^i\binom{N-1}{i}
 \binom{t+1}{\ell}
 \binom{(a-i-1)t+N-2-i}{N+j-\ell-1}.
\end{aligned}
\tag{4.10}
\]

Equations (4.9)--(4.10), with the common-endpoint subtraction, are an exact
finite expression for the Ehrhart polynomial.  Its highest-slope part is
obtained from the uniform rank-$k$ numerator by $m$ operators

\[
 T_{w,M}(f)=wx(1-x)f'+[1+(wM-1)x]f, \tag{4.11}
\]

where the parameters are positive and converge to $1,1/2,\ldots,1/m$.
The remaining terms have smaller slope and the signed size bounds recorded in
proof-dossier D3.  On the middle annulus $Y/N\le y\le N/Y$, root-of-unity
filtering gives a conjugate leading pair $2\mathcal M\cos\Phi/k$ and

\[
 \left|\frac{E}{\mathcal M}\right|
 +\left|\partial_\Phi\frac{E}{\mathcal M}\right|
 \le C_{r,q}e^{-c_{r,q}Y}+o_N(1). \tag{4.12}
\]

At the two endpoints, after the scalings $x=z/N^k$ and reciprocal
normalization, the error and its first two derivatives tend to zero on compact
sets.  The limiting numerators are obtained from the uniform kernels by the
positive operators (4.11), and their negative roots are simple.  Thus the
homotopy from the leading model to the exact one-hyperplane polynomial has no
negative multiple root in any of the three regions.  Its constant and leading
coefficients stay positive, so roots cannot escape through zero or infinity;
all roots remain simple and negative.

For $r=3$, the same finite moment formula holds, but (4.12) is replaced by
the two-saddle identity

\[
 J_{0,N}(-y^2)=\mathcal M_N(y)\cos\Phi_N(y),\qquad
 \Phi_N'(y)\ge cN/(1+y^2), \tag{4.13}
\]

and the direct and reciprocal limits are respectively products of
$1+(z\,d/dz)/j$ applied to $\cosh\sqrt z$, and the two parity kernels
$\cosh\sqrt z$ and $\sinh(\sqrt z)/\sqrt z$.  Their function and
derivative never vanish simultaneously on the negative axis, so the same
three-region homotopy closes rank three.

Finally, every additional hyperplane meets the distinguished one in at most
$r-2$ elements.  Because $q$ is fixed, packing gives only
$O_{r,q}(n^{r-3})$ such corrections; after the normalizations above their
coefficients satisfy the same strict output-box bounds as the error in
(4.12).  They therefore enter the same homotopy without changing the argument.
The full derivation of (4.9)--(4.13), including the endpoint two-jet estimates,
is proof-dossier D3.
\end{proof}


\begin{mdstatement}{Proposition 4.5 (rank-three giant-hyperplane theorem)}


There is an integer $N_{\mathrm{giant}}$ such that, for every
$n\ge N_{\mathrm{giant}}$, every connected rank-three paving matroid whose
largest effective hyperplane has size
$s\ge\lceil2n/3\rceil-1$ has only simple strictly negative $h^*$-zeros.
All legal attached hyperplanes are included.

\end{mdstatement}

\begin{proof}


Put $q=n-s$, $N=s+1$, and $m=q-1$.  The one-hyperplane numerator has a
simple negative spectrum by the exact rank-two slice representation.  Write
the complete numerator as

\[
 p_\tau=p_{\mathrm{single}}-\tau\Delta,
 \qquad 0\le\tau\le1. \tag{4.14}
\]

where $\Delta$ is the sum of all legal attached-hyperplane corrections.
If $m/N\to0$, the sublinear-complement theorem recorded as \texttt{S9} together
with its dependency revision \texttt{S9R1} in proof-dossier D3 applies uniformly; it
covers both fixed and divergent $q$.  It therefore remains to work on an
arbitrary compact interval of positive values of $m/N\subset(0,1/2]$.

On that compact interval, root-of-unity filtering writes

\[
 p_{\mathrm{single}}=J+R_{\mathrm{single}},
 \qquad J(-y^2)=\mathcal M\cos\Phi,
 \qquad \Phi'>0. \tag{4.14a}
\]

The exact parent-remainder, packing, and contour estimates give, for a fixed
sufficiently large threshold $B$,

\[
 \frac{|R_{\mathrm{single}}-\tau\Delta|}{\mathcal M}<\frac12,
 \qquad 0\le\tau\le1, \tag{4.15}
\]

at every phase-grid point with $ny^{2/3}\ge B$.  The estimate is uniform in
the compact ratio interval and in every attached family; its proof combines
the parent remainder with the full attached correction and splits the zero,
finite-positive, and far-end scales, as recorded in proof-dossier D3.  On the
remaining compact endpoint, the weighted kernel of (4.14) has initial value
one, downward variation below one third, and only upward jumps.  Consequently
the function and derivative are uniformly separated from simultaneous zero
for every $\tau\in[0,1]$.

The constant term, degree, and positive leading coefficient of (4.14) are
independent of $\tau$.  A root therefore cannot escape through zero or
infinity.  At the first parameter where the simple negative spectrum could
fail, a negative multiple root would occur.  The compact endpoint separation
excludes it below $B$, and (4.15), together with $\Phi'>0$, excludes it
above $B$.  Thus the homotopy reaches the complete polynomial at $\tau=1$
with all roots simple and negative.  A failure sequence and compactness of
$m/N\in[0,1/2]$ convert the sequential statement into one common threshold.

\end{proof}


\section{Fixed-rank exhaustion}


We now assemble the fixed-rank result.  The following three lemmas are the
uniform regime statements supplied by Sections 3--4 and the endpoint kernels.

\begin{mdstatement}{Lemma 5.1 (bounded critical excess)}


Fix $r\ge4$ and $C\ge0$.  Let

\[
 D_u=n-\left\lceil\frac nr\right\rceil,
 \qquad
 s=\begin{cases}
 \max_{H\in\mathcal H(M)}|H|,&\mathcal H(M)\ne\varnothing,\\
 r-1,&\mathcal H(M)=\varnothing.
 \end{cases}
\]

For all sufficiently large $n$, every connected rank-$r$ paving matroid
with $s\le D_u+C$ has a simple strictly negative $h^*$-spectrum.

\end{mdstatement}

\begin{proof}


There are two different endpoint problems.  If $s\le D_u-1$, put
$g=D_u-s\ge1$.  This is the infinite family $F_{r,b,g}$, not a finite
collection.  Fix sufficiently large endpoint counts $j,k$.  If the transfer
were not uniform over all $g$, choose a failure sequence.  After fixing the
finite residue $b$, either $g$ is constant on a subsequence or
$g\to\infty$; Proposition 4.1 then gives the limiting kernel
$F_{r,b,g}$ or $V_{r,b}$, respectively, on the same zero-free counting
circle.  Rouché's theorem contradicts failure.  The direct endpoint has the
same failure-sequence compactness, so both endpoint blocks have common
thresholds without taking a maximum over all $g$.

If $\delta=s-D_u\in\{0,\ldots,C\}$, write

\[
 \begin{aligned}
 b_u&=r\lceil n/r\rceil-n,\\
 \ell&=\left\lceil\frac{\delta+1-b_u}{r-1}\right\rceil,\\
 c&=\lceil n/r\rceil+\ell,\qquad
 b=b_u+r\ell,\qquad
 a=(r-1)\ell+b_u-\delta-1.
 \end{aligned} \tag{5.1}
\]

For fixed $r,C$, only finitely many $(b,a)$ occur, and the
$E_{r;b,a}$-kernels in Proposition 4.1 give common endpoint thresholds.
In the first subdomain, the phase grid (2.12) is ordered and Corollary 3.3
gives alternating strict signs.  In the bounded-excess subdomain the required
signs instead come from Proposition 3.3A.  The intrinsic degree is

\[
 n-\max\left\{\left\lceil\frac nr\right\rceil,
 \left\lceil\frac{s+1}{r-1}\right\rceil\right\}; \tag{5.2}
\]

Indeed, in the first subdomain take $c=\lceil n/r\rceil$ and
$b=rc-n$.  At dilation $c$, assign one to every coordinate and add one to
any $b$ coordinates.  Every effective hyperplane then has sum at most

\[
 s+b\le D_u-1+b=(r-1)c-1,
\]

so this is a relative interior lattice point.  In the bounded-excess domain,
use the parameters in (5.1).  The largest hyperplane is unique for large
$n$, and the exact identities give $q\ge b$; assign the $b$ extra units
to distinct coordinates outside it.  Its sum is
$s=(r-1)c-a-1<(r-1)c$.  Every other effective hyperplane has size at most
$q+r-2$, and its total is at most
$q+r-2+b=c+a+r-1<(r-1)c$.  Conversely, positivity of all coordinates and
the largest-hyperplane inequality force both lower bounds in (5.2).  Thus the
codegree and degree are exact.  The endpoint counts plus the intermediate
sign changes equal this degree.
Local Rouché disks at the endpoints and strict sign separation in the middle
also exclude multiplicity.
\end{proof}


\begin{mdstatement}{Lemma 5.2 (subcritical giant hyperplane)}


Fix $r\ge5$ and $0<p_{\max}<1/r$.  For all sufficiently large $n$, every
connected rank-$r$ paving matroid whose largest effective hyperplane has
complement size $q$ satisfying

\[
 2\le q\le p_{\max}n \tag{5.3}
\]

has a simple strictly negative $h^*$-spectrum.

\end{mdstatement}

\begin{proof}


Choose $Q=Q(r,p_{\max})$ from the uniform divergent-complement estimate.
For $Q\le q\le p_{\max}n$, distinct effective hyperplanes meet in at most
$r-2$ elements, so every other hyperplane has size at most $q+r-2$.
Proposition 4.3 treats both endpoint blocks and gives the exact degree.  The
complete middle sign is Proposition 3.4, not Proposition 3.1: the largest
hyperplane lies beyond the latter proposition's central-density domain.
The paving intersection bound places every other hyperplane in the separated
error term already included in Proposition 3.4.  Ordered first hits of the
levels between the two endpoint phases yield every remaining root.

Only the finite set $q=2,\ldots,Q-1$ remains.  For each fixed $q$, the
fixed-complement theorem, Proposition 4.4, supplies a threshold $N_q$, including all
legally attached hyperplanes.  Taking the maximum of these finitely many
thresholds and the uniform $q\ge Q$ threshold proves the lemma.  Equivalently,
any failing sequence has either a divergent-$q$ subsequence or a constant-
$q$ subsequence, and both are excluded.
\end{proof}


\begin{mdstatement}{Lemma 5.3 (critical moving giant hyperplane)}


Fix $r\ge5$.  Suppose a sequence of connected rank-$r$ paving matroids has
largest-hyperplane complement $q$ with

\[
 \frac qn\longrightarrow\frac1r,
 \qquad n-rq\longrightarrow\infty. \tag{5.4}
\]

Then its $h^*$-polynomials eventually have only simple strictly negative
zeros.

\end{mdstatement}

\begin{proof}


Fix large integers $J,k$ before sending $n\to\infty$.  At the direct
endpoint, the disk ending at $x_L=-R_{p,J}/n^r$ contains exactly $J$
simple negative roots and has boundary sign $(-1)^J$.  At the reciprocal
endpoint, after fixing the finite residue parameter, the exterior of
$x_U=-\Lambda/T_{a,k}^d$ contains exactly $k$ simple negative roots and
has boundary sign $(-1)^{D-k}$.

The four quantities

\[
 q,\qquad n-rq,\qquad nh,\qquad n/h \tag{5.5}
\]

exceed a common fixed-rank threshold throughout the interval between these
endpoint boundaries.  Here $h$ is the increasing phase-scale parameter:
along the leading branch,

\[
 K=\frac{a_0h}{dp}
 \quad\text{and}\quad
 a_s=\frac{a_0t^r}{1-p}
   =h^{r-1}\frac{-\sin(r\theta)}{\sin\theta}
\]

are strictly increasing with $\theta\in(\pi/r,\pi/(r-1))$; the saddle
equation then makes $h$ strictly increasing with $t$.
The first two quantities tend to infinity by (5.4), and the latter two are
made large by choosing $J,k$ first.  Proposition 3.4 then gives the complete
actual sign at every intermediate $U_n$-phase level.  Taking the first
point at which each level $(J+1)\pi,\ldots,(D-k-1)\pi$ is reached gives an
ordered grid without assuming a stronger global monotonicity statement.
The endpoints and grid supply $J+(D-J-k)+k=D$ distinct negative roots, so
degree exhaustion makes them all simple.  Finite residue classes are handled
by a finite subsequence maximum.
\end{proof}


\begin{mdstatement}{Lemma 5.4 (the complete rank-three class)}


There is an integer $N_3$ such that every connected rank-three paving
matroid on at least $N_3$ elements has a simple strictly negative
$h^*$-spectrum.

\end{mdstatement}

\begin{proof}


For $b\in\{0,1,2\}$ and $g\ge1$, its reciprocal kernels are

\[
 F_{b,g}(z)=\sum_{j\ge0}
 \Pr\{\operatorname{Bin}(3j+b,2/3)\le2j+b+g-1\}
 \frac{z^j}{(3j+b)!}. \tag{5.6}
\]

At
$T_{i,b}=2\pi(i+b/3)/\sqrt3$, exact negative-axis signs give
$(-1)^iF_{b,g}(-T_{i,b}^3)>0$, uniformly in $g$.  Full complex-circle
counts give exactly $i$ roots inside $|z|=T_{i,b}^3$; hence all zeros of
every $F_{b,g}$ are simple and negative.

For completeness, the sign is not a numerical extrapolation.  Writing the
kernel as the three-pole sum plus its fixed-radius remainder, the normalized
remainder $Q_R(T)$ and the positive comparison $d(T)$ satisfy

\[
 Q_R(T)<d(T)\qquad(T\ge3). \tag{5.6a}
\]

On $[3,10^6]$, (5.6a) follows from 36 adjacent rational intervals: for
each fixed rational radius, $Q_R$ is convex and $d$ is concave, so the two
endpoint inequalities certify the whole interval.  For $T\ge10^6$, the
analytic tail bound is (Q\_R/d<0.975).  Since every $T_{i,b}>3$, this proves
all signs.  The independent full-circle estimate has the same dominant
three-pole sum and a strict uniform boundary margin; the argument principle
therefore gives exactly $i$ roots.  The rational interval certificate and
the infinite-tail inequality are included in proof-dossier D3.

The direct endpoint limits are indexed by legal macroscopic spectra
$\boldsymbol\alpha=(\alpha_j)$, $\sum\alpha_j\le1$, $\max\alpha_j<1$:

\[
 F_{\boldsymbol\alpha}^{(3)}(z)=
 \sum_{\ell\ge0}\left[1-
 \sum_j\Pr\{\operatorname{Bin}(3\ell,\alpha_j)\ge2\ell+1\}\right]
 \frac{z^\ell}{(3\ell)!}. \tag{5.6b}
\]

Every function in (5.6b) has infinitely many simple negative zeros.  The proof
splits at the largest two masses.  Small spectra satisfy a direct Beta--Rouché
bound; otherwise the common signed weight has downward variation

\[
 V_-=\sum_{j\ge3}\frac{\alpha_j}{2-3\alpha_j}+1-g(\alpha_2)<\frac12, \tag{5.6c}
\]

where, if $p_b\in(2/3,1)$ is determined by
$b^2(1-b)=p_b^2(1-p_b)$,

\[
 g(b)=\frac{p_b}{3p_b-2}-\frac{b}{2-3b}. \tag{5.6c'}
\]

The three possible ranges of $\alpha_2$ reduce to the explicit
Bernstein-positive rational inequalities in proof-dossier D3.  Upward jumps,
finite total variation, and the gamma multiplier then give the asserted simple
negative spectrum.  This also pays countable spectra by dominated truncation.

Suppose a common threshold fails and choose a failing sequence.  If
$s\ge\lceil2n/3\rceil-1$ along an infinite subsequence, Proposition 4.5
excludes that subsequence; this includes both fixed and divergent complements
of the largest hyperplane.  We may therefore pass to a subsequence with
$s<\lceil2n/3\rceil-1$.  If $s/n\le a<2/3$ for some fixed $a$, packing
gives a limiting spectrum for (5.6b); the exact contour formula converges
uniformly on the direct circle, while the same size-layer estimate gives

\[
 \frac{|E|}{R^n}\le C_a(1+t^3)e^{-c_at} \tag{5.6d}
\]

on the moving middle axis.  Thus the direct roots from (5.6b), the alternating
middle grid, and the reciprocal kernel roots transfer uniformly.

The only remaining possibility is $s/n\to2/3$ from below.  Put
$c=\lceil n/3\rceil$, $D=n-c$, $b=3c-n$, and $g=D-s$.
If $g\to\infty$, the reciprocal limit is $V_{3,b}$; if $g$ is bounded,
fix the finite pair $(b,g)$ and use $F_{b,g}$.  The zero-free circles proved
above transfer the exact reciprocal block.  Independently, compactness of the
hyperplane densities gives a direct limit (5.6b), so a fixed direct block also
transfers.  Between them, the exact double-moving formula satisfies (5.6d) in
the separated part and the complementary-error estimate (3.22c) in the
bounded critical window.  Hence every integer phase has its alternating
complete-polynomial sign.  Ordered first hits give (D-j-k) middle roots;
together with the (j) direct and (k) reciprocal roots, they exhaust the
actual degree (D).  All roots are therefore simple and negative, contradicting
the chosen sequence.
\end{proof}


\begin{mdstatement}{Lemma 5.5 (the complete rank-four class)}


There is an integer $N_4$ such that every connected rank-four paving
matroid on at least $N_4$ elements has a simple strictly negative
$h^*$-spectrum.

\end{mdstatement}

\begin{proof}


Let $D_u=n-\lceil n/4\rceil$, $s=s(M)$, and $q=n-s$.  The exact
regime theorems needed here are:

\[
 \begin{array}{c|l}
 I_0&s\le D_u:\ \text{Lemma 5.1 with }C=0,\\
 I_1&D_u<s\le D_u+C:\ \text{Lemma 5.1 for fixed }C,\\
 I_2&q/n\to0:\ \text{fixed }q\text{ or }q\to\infty\text{ subcritical},\\
 I_3&q/n\to p_0\in(0,1/4):\ \text{subcritical giant profile},\\
 I_4&q/n\to1/4, n-4q\to\infty:\ \text{critical giant profile}.
 \end{array} \tag{5.7}
\]

In $I_2$, a bounded $q$ subsequence is fixed and uses the exact
fixed-complement theorem, Proposition 4.4; a divergent $q$ subsequence uses
Propositions 3.5 and 4.3.  Regimes $I_3$ and $I_4$ use the corresponding
interior and moving-critical clauses of Proposition 3.5 together with the
endpoint counts of Proposition 4.3.  In every case, the endpoint circles contain exact fixed
root counts, the appropriate phase grid has complete-polynomial signs, and
the codegree construction in Proposition 4.3 gives the actual degree; ordered
first hits therefore exhaust all roots.

It remains only to see that this finite regime list is exhaustive.  Put
$c=\lceil(s+1)/3\rceil$,
$a_n=3c-s-1\in\{0,1,2\}$, and $b=4c-n$.  If $b$ is bounded along a
failing subsequence, then $s-D_u$ is bounded and $I_1$ applies.  If
$b\to\infty$, then $n-4q=3b-4a_n-4\to\infty$, and a subsequential limit
of $q/n\in[0,1/4]$ lies in exactly one of $I_2,I_3,I_4$.  Hence no failing
rank-four sequence exists.
\end{proof}


\begin{proof}[Proof of Theorem 1.1]


First assume $M$ is connected.  Ranks three and four are Lemmas 5.4 and 5.5.

Now fix $r\ge5$, and suppose no common threshold exists.
Choose a failing sequence with $n\to\infty$.
For each member write $s=s(M)$ and
$D_u=n-\lceil n/r\rceil$, in accordance with (1.0).

If $s\le D_u$ infinitely often, Lemma 5.1 with $C=0$ is a contradiction.
Otherwise put $\delta=s-D_u>0$.  If $\delta$ is bounded on a subsequence,
Lemma 5.1 with one fixed bound $C$ is again a contradiction.  Thus we may
assume $\delta\to\infty$.  With $q=n-s$ and
$b_u=r\lceil n/r\rceil-n\in\{0,\ldots,r-1\}$,

\[
 n-rq=r\delta-b_u\longrightarrow\infty. \tag{5.8}
\]

After taking a subsequence, $q/n\to p_0\in[0,1/r]$.  If $p_0<1/r$,
choose one fixed $p_{\max}\in(p_0,1/r)$ and apply Lemma 5.2.  If
$p_0=1/r$, apply Lemma 5.3 using (5.8).  Every possible bad sequence has
been contradicted, so there is a common connected threshold.

It remains to classify disconnected paving matroids for fixed $r$ and
$n>r$.  A component containing a circuit has rank at least $r-1$, and
$2(r-1)>r$, so there is at most one such component.  Every other connected
component is a one-element coloop.  Disconnectedness and the total rank force

\[
 M=U_{r-1,n-1}\oplus U_{1,1}. \tag{5.9}
\]

Deleting the fixed coloop coordinate is an integral affine isomorphism of
base polytopes and preserves the Ehrhart polynomial.  The fixed-rank uniform
hypersimplex theorem for rank $r-1$ supplies a threshold for (5.9) when
$r-1\ge3$, by Adiprasito--Zhang [1, Theorem 3.1].
For $r=3$, put $m=n-1$.  The rank-two numerator formula [2] is
$P_m(x)=\sum_j\binom m{2j}x^j-mx$.  At $x=-\tan^2\theta$,
$0\le\theta<\pi/2$, it gives

\[
 \cos^m\theta\,P_m(-\tan^2\theta)
 =\cos(m\theta)+m\sin^2\theta\cos^{m-2}\theta.
\]

For $m\ge5$, the nonnegative second term has maximum
$2(1-2/m)^{(m-2)/2}<1$: the maximum decreases with $m$, and its
square at $m=5$ is $108/125<1$.  Hence the values at
$\theta=j\pi/m$ have sign $(-1)^j$.  If $m$ is odd, this gives
$\lfloor m/2\rfloor$ disjoint sign-changing intervals.  If $m$ is
even, the endpoint limit has sign $(-1)^{m/2}$ and supplies the last
interval.  Since $\deg P_m=\lfloor m/2\rfloor$, all roots are distinct
and strictly negative.  Thus $n\ge6$ suffices for this disconnected
rank-three subcase; simplicity is justified by this argument, not by an
unstated simplicity assertion in [2].  Taking
the maximum of the connected and disconnected thresholds proves the theorem.

\end{proof}


\section{The growing-rank root-density theorem}


We prove Theorem 1.2.  Suppress the index $\nu$, and define

\[
 J=\lceil r^7A^4\rceil,
 \qquad K=\lceil r^5A^4\rceil. \tag{6.1}
\]

Condition (1.1) implies

\[
 X\to\infty,\qquad \frac{r^4}{n}\to0,\qquad
 \log n\gg r^4A^4\gg r^3\log(2r). \tag{6.2}
\]

At the $J$-th direct-end grid point,

\[
 (n-1)t_J=rJ(1+o(1))=r^8A^4(1+o(1)). \tag{6.3}
\]

At $t_J$, the first two quantities in Proposition 3.0 are paid by
(6.2)--(6.3).  On $[t_J,t_-]$, $\kappa=(n-1)t$ increases, whereas the
logarithmic-separation and the two factors containing $t$ are weakest at the
right endpoint.  There,

\[
 \log\!\left[(C_1r)^rC_*^r t_-\log(1+\kappa)\right]
 \le-\frac{\log n}{K_0r}+O(r\log(2r)+A)\longrightarrow-\infty, \tag{6.3a}
\]

and the same estimate without $(C_1r)^r$ is stronger.  Moreover
$\log(1/t_-)=\log n/(K_0r)$ dominates
$r^2+r\log(2r)$ by (6.2).  Since
$t\log(1+(n-1)t)$ increases with $t$, these right-end estimates and the
left-end $\kappa$-estimates pay all five gates throughout
$[t_J,t_-]$.  Proposition 3.0, not merely a one-hyperplane central limit,
gives the complete $h_M^*$-sign on this ultra-small interval.  Proposition
3.1 then covers
$[t_-,1/4]$, $[1/4,4]$, and the large-$t$ interval down to the reciprocal
$K$-scale.  Corollary 3.3 gives

\[
 (-1)^jh_M^*(-t_j^r)>0,
 \qquad J\le j\le D-K. \tag{6.4}
\]

Thus the intervals between consecutive grid points contain

\[
 D-K-J \tag{6.5}
\]

distinct strictly negative zeros.

By Proposition 4.2, the reciprocal endpoint contributes $K-1$ additional
simple strictly negative zeros beyond its boundary radius.  One last zero is
forced between that boundary and the last phase-grid point.  Indeed, if
$b=r\lceil n/r\rceil-n$, the phase expansion (2.14) gives

\[
 \frac n{T_{K,b}}<t_{D-K}<\frac n{T_{K-1,b}}, \tag{6.6}
\]

while the reciprocal boundary sign is

\[
 \operatorname{sign}h_M^*(-n^r/T_{K-1,b}^r)=(-1)^{D+K-1}, \tag{6.7}
\]

opposite to the grid sign $(-1)^{D-K}=(-1)^{D+K}$.  The reciprocal block
therefore contributes $K$ zeros disjoint from (6.5).  The total is $D-J$.

It remains to justify the degree.  Put $c=\lceil n/r\rceil$ and
$b=rc-n\in\{0,\ldots,r-1\}$.  Every interior lattice point of $kP(M)$
has all coordinates at least one and coordinate sum $rk$, so necessarily
$k\ge c$.  At $k=c$, start with all $n$ coordinates equal to one and
increase any $b$ coordinates to two.  Since (1.1) implies $c>2$ eventually,
all coordinate inequalities are strict.  For every effective hyperplane,

\[
 x(H)\le |H|+b\le D-1+b=(r-1)c-1<(r-1)c, \tag{6.8}
\]

so every hyperplane inequality is also strict; lower-rank flat inequalities
follow from paving and the coordinate bounds.  Thus the codegree is exactly
$c$.  Since connectivity gives $\dim P(M)=n-1$,

\[
 \deg h_M^*=n-\left\lceil\frac nr\right\rceil=D. \tag{6.9}
\]

Finally, (1.1) and (6.2) imply $J/D\to0$.  This proves Theorem 1.2.

\hfill$\square$


\section{Scope, prior work, and limitations}


The paving Ehrhart formula and the object-level hyperplane decomposition are
prior work.  Results for rank-two sparse paving matroids and the recent
fixed-rank hypersimplex theorem are also prior work.  In particular, the
uniform subfamily of Theorem 1.1 is not claimed as new.

Adiprasito and Zhang [1, Theorem 3.1] prove that, for each fixed $k\ge3$,
the ordinary hypersimplex $\Delta(k,n)$ has a numerator with only distinct
negative real zeros for all sufficiently large $n$.  Their Theorem 1.1
gives negative real-rootedness for rank three and $n\ge6$, with simplicity
asserted for $n\ge32$.  Theorem 1.1 here extends the object scope of their
eventual fixed-rank theorem from uniform to all paving matroids; it does not
improve their small-order rank-three range.

They also remark that their threshold can be quantified as
$N(k)=k^{O(k^2)}$.  With that growth bound, the uniform specialization in
the regime (1.1) is already covered by their stronger full-root conclusion.
Rank growth alone is therefore not the new part of Theorem 1.2: its increment
is the nonuniform paving range and the explicit root-density window, with the
stated hyperplane restriction.  Theorem 1.2 is not a strengthening of their
full-root theorem on uniform matroids.

The residue entire functions in [1, Section 3.3] and the phase-grid,
endpoint-transfer, and degree-exhaustion architecture in its Sections
3.4--3.5 are prior analytic mechanisms used here.  The full text of version
v1 was checked; it contains no theorem covering all nonuniform paving
matroids or the nonuniform conclusion of Theorem 1.2.  This source comparison
does not establish global priority over every other or unindexed work.

Theorem 1.1 supplies no explicit $N_r$ and no threshold uniform in $r$.
Theorem 1.2 assumes connectivity, (1.1), and the effective-hyperplane size
bound; it leaves $o(D)$ zeros unclassified.  Neither theorem resolves a
general real-rootedness conjecture for arbitrary matroid base polytopes.

\section{Reproducibility}


The frozen prewrite package contains 216 source, proof, and certificate files.
Their ordered path-and-SHA-256 manifest is

loops/PAVING-HSTAR-PUBLICATION-PREP-0001/A1-{\cjkfont 证\allowbreak{}明\allowbreak{}依\allowbreak{}赖\allowbreak{}闭\allowbreak{}包}.sha256.

The finite certificates in the package check only explicitly stated rational
inequalities and finite continuous covers.  They do not replace the analytic
proofs above.  The manuscript and any public artifact must be rebuilt if a
file in the frozen manifest changes.

For auditability, the principal manuscript assertions map to the frozen proof
groups as follows; the manifest fixes the exact bytes of every named file and
its bound revisions.

\begin{itemize}
\item Equations (2.1)--(2.9): \texttt{S10}, \texttt{S10R1}, and \texttt{S11} in
\texttt{PAVING-\allowbreak{}HSTAR-\allowbreak{}HIGHER-\allowbreak{}RANK-\allowbreak{}CRITICAL-\allowbreak{}KERNEL-\allowbreak{}GATE-\allowbreak{}0001}.
\item Proposition 3.0: \texttt{S17} and \texttt{S20-\allowbreak{}S17{\cjkfont 门\allowbreak{}下\allowbreak{}的\allowbreak{}完\allowbreak{}整\allowbreak{}原\allowbreak{}端\allowbreak{}保\allowbreak{}号}} in
\texttt{PAVING-\allowbreak{}HSTAR-\allowbreak{}CROSS-\allowbreak{}RANK-\allowbreak{}FACT-\allowbreak{}GATE-\allowbreak{}0001}, together with that gate's \texttt{S13}.
\item Propositions 3.1--3.2: \texttt{S18}, \texttt{S14}, and \texttt{S15} in
\texttt{PAVING-\allowbreak{}HSTAR-\allowbreak{}CROSS-\allowbreak{}RANK-\allowbreak{}FACT-\allowbreak{}GATE-\allowbreak{}0001}, plus the full low-Euler estimate
in \texttt{S11}.
\item Proposition 3.3A: \texttt{S59} and \texttt{S91} in
\texttt{PAVING-\allowbreak{}HSTAR-\allowbreak{}HIGHER-\allowbreak{}RANK-\allowbreak{}CRITICAL-\allowbreak{}KERNEL-\allowbreak{}GATE-\allowbreak{}0001}.
\item Proposition 3.4: the full geometry, moving, actual-correction, endpoint, and
critical-moving groups listed in proof-dossier D4.
\item Proposition 3.5: \texttt{S50}--\texttt{S56} in
\texttt{PAVING-\allowbreak{}HSTAR-\allowbreak{}HIGHER-\allowbreak{}RANK-\allowbreak{}CRITICAL-\allowbreak{}KERNEL-\allowbreak{}GATE-\allowbreak{}0001}, as itemized in
proof-dossier D5.
\item Proposition 4.0: the complete macro-spectrum theorem in
\texttt{PAVING-\allowbreak{}HSTAR-\allowbreak{}TWO-\allowbreak{}MACRO-\allowbreak{}GATE-\allowbreak{}0001/\allowbreak{}S3}, with its frozen analytic dependencies.
\item Proposition 4.1 and Lemma 5.1: \texttt{S90} and \texttt{S91} in
\texttt{PAVING-\allowbreak{}HSTAR-\allowbreak{}HIGHER-\allowbreak{}RANK-\allowbreak{}CRITICAL-\allowbreak{}KERNEL-\allowbreak{}GATE-\allowbreak{}0001}.
\item Proposition 4.3: the Wright, reciprocal, phase-joining, and critical endpoint
groups in proof-dossier D4.
\item Proposition 4.4: \texttt{S1} and \texttt{S2} in
\texttt{PAVING-\allowbreak{}HSTAR-\allowbreak{}FIXED-\allowbreak{}OUTSIDE-\allowbreak{}GATE-\allowbreak{}0001}.
\item Proposition 4.5: \texttt{S20-\allowbreak{}{\cjkfont 全\allowbreak{}部\allowbreak{}巨}H{\cjkfont 依\allowbreak{}赖\allowbreak{}整\allowbreak{}类\allowbreak{}统\allowbreak{}一\allowbreak{}实\allowbreak{}际\allowbreak{}全\allowbreak{}轴\allowbreak{}候\allowbreak{}选}} in
\texttt{PAVING-\allowbreak{}HSTAR-\allowbreak{}GIANT-\allowbreak{}H-\allowbreak{}DEPENDENCIES-\allowbreak{}GATE-\allowbreak{}0001}, together with the giant-
hyperplane moving and endpoint groups itemized in proof-dossier D4.
\item Lemma 5.2: \texttt{S8} in \texttt{PAVING-\allowbreak{}HSTAR-\allowbreak{}GENERAL-\allowbreak{}RANK-\allowbreak{}GIANT-\allowbreak{}ENDPOINT-\allowbreak{}GATE-\allowbreak{}0001}
and the fixed-complement theorem in \texttt{PAVING-\allowbreak{}HSTAR-\allowbreak{}FIXED-\allowbreak{}OUTSIDE-\allowbreak{}GATE-\allowbreak{}0001}.
\item Lemma 5.3: \texttt{S4} in
\texttt{PAVING-\allowbreak{}HSTAR-\allowbreak{}GENERAL-\allowbreak{}RANK-\allowbreak{}GIANT-\allowbreak{}CRITICAL-\allowbreak{}MOVING-\allowbreak{}GATE-\allowbreak{}0001}.
\item The rank-three and rank-four joins in Theorem 1.1: \texttt{S22}, \texttt{S22R1}, \texttt{S23},
and \texttt{S24} in \texttt{PAVING-\allowbreak{}HSTAR-\allowbreak{}MOVING-\allowbreak{}CRITICAL-\allowbreak{}GATE-\allowbreak{}0001}, and \texttt{S57}
with its audited interfaces in
\texttt{PAVING-\allowbreak{}HSTAR-\allowbreak{}HIGHER-\allowbreak{}RANK-\allowbreak{}CRITICAL-\allowbreak{}KERNEL-\allowbreak{}GATE-\allowbreak{}0001}.
\item Proposition 4.2 and Section 6: \texttt{S15}, \texttt{S19}, and \texttt{S20} in
\texttt{PAVING-\allowbreak{}HSTAR-\allowbreak{}CROSS-\allowbreak{}RANK-\allowbreak{}FACT-\allowbreak{}GATE-\allowbreak{}0001}.
\end{itemize}

These are proof dependencies, not separate novelty claims.  Novelty and source
coverage are recorded outside the 216-file proof closure in the publication
preparation package.

\section*{References}


\begin{enumerate}
\item K. Adiprasito and M. Zhang, \emph{Eventual unimodality of the hypersimplex}, HAL preprint hal-05747878v1, submitted 12 September 2026; Theorems 1.1 and 3.1 and Sections 3.3--3.5.
\item L. Ferroni, K. Jochemko, and B. Schröter, \emph{Ehrhart polynomials of rank two matroids}, arXiv:2106.08183v3.
\item H. Fu, \emph{Real-rootedness and ultra log-concavity of rank-two matroid Ehrhart h*-polynomials}, arXiv:2609.37439v1.
\item D. Hanely, J. Martin, J. McGinnis, D. Miyata, G. Nasr, A. Vindas-Meléndez, and M. Yin, \emph{Ehrhart Theory of Paving and Panhandle Matroids}, arXiv:2201.12442v3.
\item S. Li, \emph{Magic positivity of Snapper polynomials for matroids}, arXiv:2609.24917v1.
\end{enumerate}

\section*{AI-assisted research disclosure}


The research, proof development, checking, and writing were conducted with
substantial AI assistance.  Carptopus is the human author and assumes
responsibility for the mathematical claims and any released version.
\end{document}
